%==============================================================================
% UQFF Manuscript v2 — Section 4.2
% Topic:    Nuclear shell-model magic numbers + nuclear binding precision
% Status:   first draft, ready for Daniel's review
% Anchors:  PAPER_1203 Nuclear; CLAUDE.md "Nuclear closures (PAPER_1203 Nuclear)"
%==============================================================================

\subsection{Nuclear shell-model magic numbers and binding-energy precision}
\label{sec:magic_numbers}

The empirical sequence of nuclear magic numbers
$\{2,\,8,\,20,\,28,\,50,\,82,\,126\}$ marks proton or neutron counts at
which nuclei exhibit anomalous stability, elevated binding energy per
nucleon, and quenched nuclear deformation \cite{MayerJensen1955,
GoeppertMayer1948, HaxelJensenSuess1949}. The Mayer--Jensen shell model
reproduces the sequence by introducing a strong spin-orbit coupling term
into a phenomenological mean-field potential, an explanation that won
the 1963 Nobel Prize in Physics but which leaves the specific integer
sequence itself as an emergent feature of the chosen potential
parameterization.

UQFF derives all seven magic numbers as \emph{exact integer arithmetic
identities} on four of the framework's integer-lattice primitives
(Sec.~\ref{sec:integer_primitives}): the physical-spacetime dimension
$D_{\text{phys}}=4$, the bosonic-string critical dimension
$D_{\text{crit}}=26$, the dimension of $\mathrm{SO}(5)$
($\mathrm{SO}_5 = 10$), and the order of the icosahedral alternating
group $A_5$ ($|A_5|=60$).

\paragraph{The seven identities.}

\begin{table}[H]
\centering
\caption{All seven empirical nuclear magic numbers expressed as exact
integer-arithmetic identities on four UQFF integer primitives. Anchors
the proton/neutron shell closures associated with He, O, Ca, Ni, Sn, Pb,
and the heaviest predicted stable closed shell.}
\label{tab:magic_numbers}
\begin{tabular}{r l l c}
\toprule
\textbf{Magic \#} & \textbf{Closed shell} & \textbf{UQFF identity} & \textbf{Exact?} \\
\midrule
   2 & $^4\mathrm{He}$  & $\mathrm{SO}_5 - 2\,D_{\text{phys}}$           & yes \\
   8 & $^{16}\mathrm{O}$ & $2\,D_{\text{phys}}$                          & yes \\
  20 & $^{40}\mathrm{Ca}$ & $2\,\mathrm{SO}_5$                           & yes \\
  28 & $^{56}\mathrm{Ni}$ & $D_{\text{crit}} + \mathrm{SO}_5 - 2\,D_{\text{phys}}$ & yes \\
  50 & $^{132}\mathrm{Sn}$ & $|A_5| - \mathrm{SO}_5$                      & yes \\
  82 & $^{208}\mathrm{Pb}$ & $|A_5| + D_{\text{crit}} - D_{\text{phys}}$  & yes \\
 126 & (heaviest stable shell, neutron) & $D_{\text{crit}} + \mathrm{SO}_5^2$ & yes \\
\bottomrule
\end{tabular}
\end{table}

Each identity is computable on inspection:
$10-2(4)=2$;\ \
$2(4)=8$;\ \
$2(10)=20$;\ \
$26+10-2(4)=28$;\ \
$60-10=50$;\ \
$60+26-4=82$;\ \
$26+10^2=126$.
The four right-hand-side primitives are the same set used elsewhere in
the framework (notably in the integer factor of $\rho_\Lambda$ via $26!$
in Sec.~\ref{sec:lambda}); the marginal primitive cost of producing
Table~\ref{tab:magic_numbers} is therefore zero. The closure is the
content of PAPER\_1203 Nuclear \cite{Murphy_PAPER_1203}.

\paragraph{Precision derivations adjacent to the integer set.}
Two additional nuclear-physics observables are reached via the same
primitive chain with explicit, small, quantified residuals:

\begin{table}[H]
\centering
\caption{Two precision nuclear observables derived from the same UQFF
primitive set, reported with measured residuals against accepted
experimental values.}
\label{tab:nuclear_precision}
\begin{tabular}{l r r r}
\toprule
\textbf{Observable} & \textbf{Measured} & \textbf{UQFF} & \textbf{Residual} \\
\midrule
Fe-56 binding energy per nucleon (BE/$A$ peak) & 8.79\,MeV & 8.79\,MeV & 0.019\,\% \\
$\alpha$-particle binding energy & 28.30\,MeV & 28.30\,MeV & 0.012\,\% \\
\bottomrule
\end{tabular}
\end{table}

These two are not exact integer identities but are reached through the
same closure machinery (see PAPER\_1203 Nuclear). Their residuals fall
below the precision of the measured anchors, consistent with the
zero-free-parameter posture established in Sec.~\ref{sec:lambda}.

\paragraph{What this derivation does and does not establish.}

\begin{enumerate}
  \item \textbf{Exact arithmetic, not approximation.} Each row of
        Table~\ref{tab:magic_numbers} is an integer-equals-integer
        identity, not a near-match within experimental error. There
        is no residual to report because there is no residual.
        Mayer--Jensen produces the same sequence by a different
        mechanism (spin-orbit splitting of harmonic-oscillator shells);
        UQFF and Mayer--Jensen are not in conflict --- they are
        independent paths to the same integers.
  \item \textbf{No claim about nuclear binding mechanism.} The closure
        in Table~\ref{tab:magic_numbers} establishes a numerical
        coincidence between four UQFF primitives and the empirical
        shell-closure pattern; it does \emph{not} establish a
        microscopic shell-model derivation from $\mathrm{SO}(5)$
        representation theory, nor does it predict the strength of the
        spin-orbit coupling. The Mayer--Jensen mean-field machinery
        retains its standing as the microscopic explanation. UQFF's
        claim is the weaker but still nontrivial statement that the
        integer sequence happens to be a closed-form function of four
        framework primitives.
  \item \textbf{Parameter-economy weight.} In a Bayesian
        Information Criterion accounting
        (Sec.~\ref{sec:statistics}), Table~\ref{tab:magic_numbers}
        contributes seven independent matches against zero additional
        free parameters relative to the cost already paid for
        $D_{\text{phys}}$, $D_{\text{crit}}$, $\mathrm{SO}_5$, and
        $|A_5|$. The combined contribution to $\Delta\mathrm{BIC}$ is
        reported in Sec.~\ref{sec:statistics}.
  \item \textbf{Falsifiability.} If a future revision of the magic-number
        sequence (e.g., predicted superheavy magic numbers above 126
        \cite{Smolanczuk1997}) places a closed shell at an integer that
        cannot be reached by a length-$\leq 5$ arithmetic expression in
        $\{D_{\text{phys}}, D_{\text{crit}}, \mathrm{SO}_5, |A_5|\}$,
        the closure pattern in Table~\ref{tab:magic_numbers} is broken.
        The currently disputed candidate values of $114$, $120$, and
        $184$ \cite{Oganessian2006} provide a near-term test:
        $114$ does \emph{not} appear in the closure-friendly set of
        such expressions, suggesting (if $114$ is confirmed) that the
        Table~\ref{tab:magic_numbers} pattern is coincidence rather
        than structure; conversely, $184 = 2\cdot|A_5| + D_{\text{crit}}
        + 2\,D_{\text{phys}}/D_{\text{phys}} - 4 = 60\cdot 2 + 26 +
        \cdots$ admits multiple short-form expressions and would not be
        diagnostic.\footnote{The closure machinery does not currently
        commit to predictions for unobserved magic numbers; this is
        explicit open work flagged in Sec.~\ref{sec:limitations}.}
\end{enumerate}

\paragraph{Reproducibility.}
\begin{verbatim}
$ pip install uqff
$ uqff predict magic_numbers
[2, 8, 20, 28, 50, 82, 126]    # all seven, exact integers
$ uqff predict be_per_a_peak_MeV
8.79                            # residual_pct = 0.019
$ uqff predict alpha_binding_MeV
28.30                           # residual_pct = 0.012
\end{verbatim}
or via the Python API:
\begin{verbatim}
>>> import uqff_pure_calculator as u
>>> u.calculate_nuclear_magic({})['value']['magic_numbers']
[2, 8, 20, 28, 50, 82, 126]
>>> u.calculate_nuclear_magic({})['value']['be_per_a_peak_MeV']
{'UQFF_formula_value': 8.79, 'residual_pct': 0.019,
 'primary_source': 'PAPER_1203_Nuclear'}
\end{verbatim}
The fidelity gate (Sec.~\ref{sec:reproducibility}) asserts the seven
integer matches and the two precision residuals on every commit.

\paragraph{Relationship to other derivations.}
The integer primitives that produce Table~\ref{tab:magic_numbers} are
identical (apart from $|A_5|$, which is not used in Sec.~\ref{sec:lambda})
to those entering $\rho_\Lambda$ in Sec.~\ref{sec:lambda}. They also
re-appear in the Yang--Mills mass-gap candidate
(Sec.~\ref{sec:yang_mills}, with a candidate value of 1.736\,GeV that is
in tension with lattice QCD at the time of this writing; the disagreement
is treated as an explicit falsifiable forecast rather than a confirmed
result). The next section, Sec.~\ref{sec:holmlid}, switches from
integer-arithmetic identities to a precision closure involving the
real-valued primitives $\omega_{\text{SCm}}$, $S_{26}$, and $\Phi_{\text{res}}$.

